% !TeX root = ../概率论与数理统计笔记.tex
% 由主文档通过 \input 载入；请编译主文档。
\setcounter{chapter}{3} % 保留原书章号。
\chapter{多维随机变量及其分布}
\label{chap:ps-multivariate}

本章研究多个随机量之间的联合关系，回答“分别知道各自的分布是否足够”的问题。
先修内容为第\ref{chap:ps-univariate}章及二重积分；学习目标是掌握联合、边缘、条件分布与独立性之间的联系。

\section{随机向量与联合分布}
\label{sec:ps-joint-distributions}
定义\emph{随机向量}（random vector）与\emph{联合分布}（joint distribution）。
以二维情形建立联合分布函数、联合质量函数与联合密度，明确分布的支撑区域。
训练通过事件区域求和或积分，避免只写积分核而遗漏区域约束。

\section{边缘分布与条件分布}
\label{sec:ps-marginal-conditional}
由联合分布推出\emph{边缘分布}（marginal distribution）与\emph{条件分布}（conditional distribution）。
离散情形要求被条件化取值具有正概率；具有联合密度时，在边缘密度为正处给出条件密度公式。
解释连续条件分布并非把两个零概率事件直接相除，并说明条件分布在零概率取值上的版本可以不唯一。

\section{随机变量的独立性}
\label{sec:ps-variable-independence}
用联合分布对边缘分布的分解定义随机变量的独立性，给出分布函数及质量函数、密度的相应判据。
说明密度分解按几乎处处理解，检查乘积分解时必须包含支撑区域。
证明独立随机变量分别作可测函数变换后仍独立，并区分两两独立与相互独立。

\section{常见多维分布}
\label{sec:ps-multivariate-models}
介绍多项分布（multinomial distribution），由互斥类别的独立重复试验导出其质量函数。
介绍二维正态分布（bivariate normal distribution）的参数与密度结构；一般多元正态模型在矩阵工具齐备后补充。
指出两个正态边缘分布不保证联合正态，也不保证独立；协方差判据在第\ref{chap:ps-moments}章建立。

\section{例题、习题与学习检查}
规划三角形区域上的均匀分布、带约束的联合质量表、给定总次数的类别计数等例题。
习题要求从同一联合模型计算边缘和条件概率，并构造边缘相同而联合关系不同的两个模型。
学习检查包括区域边界、归一化、条件分母及独立性依据。






